\section{Introduction}

The current version of the Universal Dependencies (UD) collections of syntactic treebanks allows us to  evaluate (a number of)  universal properties on a set of 180 languages \citep{demarneffe2021universal}. In this paper, we study the (typological) validity of  the so-called Menzerath-Altman law (MAL) on the nature of complexity in linguistic units in terms of the relationship between the size of a given unit and the size of its components in any given domain \citep{menzerath1928einige, menzerath1954architektonik, altmann1980prolegomena}. Roughly speaking, the bigger the number of components of a given unit, the smaller their respective length. In the verbal domain, for instance, this implies that the length of verbal arguments decreases as the valency or number of verbal dependents increases \citep{macutek2017menzerath}. Our aim is also to propose a reliable metric to evaluate this hypothesis in the (unbalanced) UD collection with annotated corpora of significantly different sizes - including the so-called LOL languages, with millions of tokens, but also a number of corpora from less-studied languages with less than a thousand tokens. Furthermore, in order to be able to make typologically relevant claims based on an ad hoc language sample which is not construed for such purposes in the first place, we compare a number of typologically valid subsets of our sample to ensure that our conclusions are replicable beyond the sample as a whole and are not mere artefacts of sampling bias. Another important contribution of our study is that we looked at preverbal and postverbal domains separately. That is, we calculated two additional sets of metrics in addition to our general (or combined) MAL metrics, respectively LMAL and RMAL (for left/preverbal and right/postverbal arguments). 
Our results show that MAL is a clear typologically wide-spread preference across the languages of the world, but it is not an absolute universal principle. Importantly, we find that while MAL is a strong statistical cross-linguistic preference, there are not only languages in which this preference is trivial, but also languages which display the opposite tendency.  Furthermore, our results show a clear difference between the pre- and postverbal domains. On the one hand, the MAL preference is stronger in the postverbal domain and, on the other hand, the anti MAL preference is stronger in the preverbal domain. Importantly, our findings show that in the postverbal domain VO languages are more likely to show a strong RMAL preference and inversely OV languages are more likely to show a strong LMAL preference in the preverbal domain. 

In the following, we first briefly present MAL and the previous studies on which we are building (Section \ref{MAL}). We next present our data extraction based on the UD annotation (Section \ref{UD}) and the different metrics we use in this study and motivate our methodological choices (Section \ref{metrics}). We then present our results and discuss our findings (Section \ref{results}).